% GENERATED FILE. Edit canonical lessons, then run npm run generate:books.
% canonical-source-hash: fnv1a64:4d50a9c18a430603
% canonical-lessons: LA-C189-cognatus, LA-C189-avunculus, LA-C189-matertera, LA-C189-veritas, LA-C189-mendacium

\chapter{Kin, Uncle, Aunt, Truth, Lie}
\label{ch:la-blood-relative-mother-s-brother-mother-s-sister-truth-lie}
\hlchaptermodality{189}

\begin{chapteropening}
\textbf{By the end of this chapter:} \emph{I can say a blood relative, a mother's brother, a mother's sister, truth and a lie.}

Complete the last lesson of chapter 189: I can say a blood relative, a mother's brother, a mother's sister, truth and a lie.
\end{chapteropening}

\section[cognātus, cognātī]{cognātus, cognātī — a blood relative}

\label{lesson:LA-C189-cognatus}



\begin{quote}
Before the new one: say the Latin for a farm manager, then the Latin for a tenant farmer.
\end{quote}

\subsection*{You'll want to know: cognātus, cognātī}
\textbf{cognātus, cognātī} — \textquotedblleft{}a blood relative\textquotedblright{}.

\begin{cousinweb}
\textbf{cognātus} is built on \emph{nātus}, \textquotedblleft{}born\textquotedblright{} (Lewis \& Short), older \emph{gnātus}, with \emph{con-}, \textquotedblleft{}together\textquotedblright{}: born into the same family.
\end{cousinweb}

\begin{grammarlens}[title={What you've built}]
One more word for this chapter.
\end{grammarlens}

\subsection*{Your turn}
\begin{itemize}
\raggedright
  \item \emph{Say it:} \emph{cognātus, cognātī}
  \item \emph{Say it:} \emph{cognātus, cognātī}, once more
  \item \emph{Say it:} say \emph{cognātus, cognātī}
  \item \emph{Recall:} say the Latin for a farm manager, then the Latin for a tenant farmer, then say \emph{cognātus, cognātī} again
\end{itemize}

\begin{tcolorbox}[breakable,colback=teal!4,colframe=teal!35!black,title={Before you move on}]
What does cognātus, cognātī mean? (\textquotedblleft{}A blood relative\textquotedblright{}.) Say it once more.
\end{tcolorbox}
\section[avunculus, avunculī]{avunculus, avunculī — a mother's brother, an uncle}

\label{lesson:LA-C189-avunculus}



\begin{quote}
Before the new one: say the Latin for a tenant farmer, then the Latin for a blood relative.
\end{quote}

\subsection*{You'll want to know: avunculus, avunculī}
\textbf{avunculus, avunculī} — \textquotedblleft{}a mother's brother, an uncle\textquotedblright{}.

\begin{cousinweb}
\textbf{avunculus} is a diminutive built on \emph{avus}, \textquotedblleft{}a grandfather\textquotedblright{} (Lewis \& Short): \textquotedblleft{}little grandfather\textquotedblright{}, an old name for the mother's brother. English \emph{uncle} comes from it through French.
\end{cousinweb}

\begin{grammarlens}[title={What you've built}]
One more word for this chapter.
\end{grammarlens}

\subsection*{Your turn}
\begin{itemize}
\raggedright
  \item \emph{Say it:} \emph{avunculus, avunculī}
  \item \emph{Say it:} \emph{avunculus, avunculī}, once more
  \item \emph{Say it:} say \emph{avunculus, avunculī}
  \item \emph{Recall:} say the Latin for a tenant farmer, then the Latin for a blood relative, then say \emph{avunculus, avunculī} again
\end{itemize}

\begin{tcolorbox}[breakable,colback=teal!4,colframe=teal!35!black,title={Before you move on}]
What does avunculus, avunculī mean? (\textquotedblleft{}A mother's brother, an uncle\textquotedblright{}.) Say it once more.
\end{tcolorbox}
\section[mātertera, māterterae]{mātertera, māterterae — a mother's sister, an aunt}

\label{lesson:LA-C189-matertera}



\begin{quote}
Before the new one: say the Latin for a blood relative, then the Latin for a mother's brother.
\end{quote}

\subsection*{You'll want to know: mātertera, māterterae}
\textbf{mātertera, māterterae} — \textquotedblleft{}a mother's sister, an aunt\textquotedblright{}.

\begin{cousinweb}
\textbf{mātertera} is built on \emph{māter}, \textquotedblleft{}mother\textquotedblright{} (Lewis \& Short), with an ending like the one in \emph{altera}, \textquotedblleft{}the other of two\textquotedblright{}: a second mother.
\end{cousinweb}

\begin{grammarlens}[title={What you've built}]
One more word for this chapter.
\end{grammarlens}

\subsection*{Your turn}
\begin{itemize}
\raggedright
  \item \emph{Say it:} \emph{mātertera, māterterae}
  \item \emph{Say it:} \emph{mātertera, māterterae}, once more
  \item \emph{Say it:} say \emph{mātertera, māterterae}
  \item \emph{Recall:} say the Latin for a blood relative, then the Latin for a mother's brother, then say \emph{mātertera, māterterae} again
\end{itemize}

\begin{tcolorbox}[breakable,colback=teal!4,colframe=teal!35!black,title={Before you move on}]
What does mātertera, māterterae mean? (\textquotedblleft{}A mother's sister, an aunt\textquotedblright{}.) Say it once more.
\end{tcolorbox}
\section[vēritās, vēritātis]{vēritās, vēritātis — truth}

\label{lesson:LA-C189-veritas}



\begin{quote}
Before the new one: say the Latin for a mother's brother, then the Latin for a mother's sister.
\end{quote}

\subsection*{You'll want to know: vēritās, vēritātis}
\textbf{vēritās, vēritātis} — \textquotedblleft{}truth\textquotedblright{}.

\begin{cousinweb}
\textbf{vēritās} is built on \emph{vērus}, \textquotedblleft{}true\textquotedblright{} (Lewis \& Short), with \emph{-tās}, an ending for a quality. \emph{vērus} is usually traced to Proto-Indo-European \emph{*weh1ros}, \textquotedblleft{}true\textquotedblright{} (EtymDB 2.1).
\end{cousinweb}

\begin{grammarlens}[title={What you've built}]
One more word for this chapter.
\end{grammarlens}

\subsection*{Your turn}
\begin{itemize}
\raggedright
  \item \emph{Say it:} \emph{vēritās, vēritātis}
  \item \emph{Say it:} \emph{vēritās, vēritātis}, once more
  \item \emph{Say it:} say \emph{vēritās, vēritātis}
  \item \emph{Recall:} say the Latin for a mother's brother, then the Latin for a mother's sister, then say \emph{vēritās, vēritātis} again
\end{itemize}

\begin{tcolorbox}[breakable,colback=teal!4,colframe=teal!35!black,title={Before you move on}]
What does vēritās, vēritātis mean? (\textquotedblleft{}Truth\textquotedblright{}.) Say it once more.
\end{tcolorbox}
\section[mendācium, mendāciī]{mendācium, mendāciī — a lie}

\label{lesson:LA-C189-mendacium}



\begin{quote}
Before the new one: say the Latin for a mother's sister, then the Latin for truth.
\end{quote}

\subsection*{You'll want to know: mendācium, mendāciī}
\textbf{mendācium, mendāciī} — \textquotedblleft{}a lie\textquotedblright{}.

\begin{cousinweb}
\textbf{mendācium} is built on \emph{mendāx}, \textquotedblleft{}lying, a liar\textquotedblright{} (Lewis \& Short). Where \emph{mendāx} comes from is uncertain; a link with \emph{mendum}, \textquotedblleft{}a fault\textquotedblright{}, is often suggested.
\end{cousinweb}

\begin{grammarlens}[title={What you've built}]
One more word for this chapter.
\end{grammarlens}

\subsection*{Your turn}
\begin{itemize}
\raggedright
  \item \emph{Say it:} \emph{mendācium, mendāciī}
  \item \emph{Say it:} \emph{mendācium, mendāciī}, once more
  \item \emph{Say it:} say \emph{mendācium, mendāciī}
  \item \emph{Recall:} say the Latin for a mother's sister, then the Latin for truth, then say \emph{mendācium, mendāciī} again
\end{itemize}

\begin{tcolorbox}[breakable,colback=teal!4,colframe=teal!35!black,title={Before you move on}]
What does mendācium, mendāciī mean? (\textquotedblleft{}A lie\textquotedblright{}.) Say it once more.
\end{tcolorbox}
